\section*{Appendix E: The Automation Evaluation Card}

\begin{table*}[t]
\centering
\small
\caption{The Ten-Point Automation and AI Workflow Evaluation Card}
\begin{tabularx}{\textwidth}{p{0.36\textwidth}X}
\toprule
\textbf{Evaluation Question} & \textbf{Empirical Evidence Required} \\
\midrule
What specific task is automated? & Exact input schema, output format, human owner, and baseline manual flow \\
How often does the task occur? & Observed job volume and variability, not an aspirational forecast \\
What is the baseline manual standard? & Documented grading rubric and observed human time distribution \\
What does the tool actually replace? & Direct comparison of end-to-end process with and without the tool \\
What human review is still necessary? & Sample audit frequency, confidence thresholds, and escalation rules \\
What are the comprehensive tool costs? & Token inference, software licenses, QA, monitoring, and maintenance \\
What happens when the output is wrong? & Defect consequence analysis, student remedy, and rollback procedures \\
What is the net outcome effect? & Measured accuracy, learning gain, turnaround, and student effort \\
What privacy and IP rights are needed? & Explicit learner consent, data retention limits, and data minimization \\
When should the tool be retired? & Performance drift triggers, cost ceilings, and quality degradation floors \\
\bottomrule
\end{tabularx}
\end{table*}

\subsection*{E1. A Test for ``AI Saves 70\% of Teacher Time''}

Ask which teacher tasks were timed, what percentage of the teacher's \emph{whole week} they represent, whether the comparison includes checking AI output, whether student needs and task difficulty are comparable, and whether independent learning outcomes were preserved. A tool that shortens comment drafting by 70\% may save much less than 70\% of a teacher's total workload: and the effect could reverse if errors demand extra correction.

\section*{Appendix F: The Educational Quality-and-Cost Frontier}

Treat these as \textbf{separate observed dimensions}, not ingredients of an invented master score.

\begin{table*}[t]
\centering
\small
\caption{The Eight Dimensions of the Educational Quality-and-Cost Frontier}
\begin{tabularx}{\textwidth}{p{0.18\textwidth}p{0.38\textwidth}X}
\toprule
\textbf{Dimension} & \textbf{Illustrative Observable Measure} & \textbf{Common False Inference} \\
\midrule
Access & Qualified learners able to participate & More enrollments always mean more access to value \\
Exposure & Sessions and practice opportunities & Attendance proves learning \\
Reliability & Timely reviews, service uptime, resolved incidents & Fast responses are always good responses \\
Competence & Independent unseen task under clear conditions & A polished submission proves authorship \\
Durability & Delayed or later-context performance & High immediate test scores prove retention \\
Cost & Total cash and normalized skilled-hour cost & Founder time is free \\
Customer sacrifice & Time, expense, effort, switching, lost opportunities & Low fee means low total cost \\
Sustainability & Repeatable contribution and cash reserves & A profitable-looking pilot proves scalable economics \\
\bottomrule
\end{tabularx}
\end{table*}

\paragraph*{Interpretation Protocol.} Before changing a process, state the performance standard you refuse to weaken and the specific costs you expect to reduce. After changing it, review whether quality has improved, held steady, worsened, or remains unknown. Report denominator changes and differences in learner readiness. Do not claim a frontier shift simply because you changed the scoring rubric.

\section*{Appendix G: A Hypothetical Eight-Week Delivery-Learning Sprint}

This is a planning exercise for a capital-constrained service founder, not a proven optimal schedule. It is more realistic to complete 1 bounded loop well than to create a large ``learning platform'' that has never been used.

\paragraph*{Week 1: Write the promise and full process.} Define a specific learner or client, the outcome, the limits, and each task from sales to assessment. Interview likely buyers about their expectations and time commitments; do not substitute praise for actual demand.

\paragraph*{Week 2: Measure the work in miniature.} Run a representative diagnostic, explanation, practice task, review, and assessment using authentic non-confidential material. Time each activity. Identify the likely bottleneck.

\paragraph*{Week 3: Design the minimum repeatable system.} Use existing platforms, a clear rubric, a version-controlled repository, a response-time policy, and simple financial tracking. Do not build unnecessary software.

\paragraph*{Week 4: Obtain real commitments.} When the offer can be delivered lawfully and reliably, present clear price, scope, dates, prerequisites, and cancellation terms to qualified prospects. Confirm paid or properly authorized pilot demand; record objections and workload.

\paragraph*{Week 5: Observe production without hiding errors.} Deliver the 1st part. Record queues, exceptional cases, learner attempts, support time, and incidents. Identify whether support needs arise from meaningful learning difficulty or avoidable process friction.

\paragraph*{Week 6: Improve the highest-cost recurring failure.} Choose 1 bottleneck: a broken setup process, unclear assignment, repeated misconception, or reviewer backlog. Introduce 1 targeted intervention and compare the end-to-end result.

\paragraph*{Week 7: Test independence and full costs.} Use an unfamiliar assessment task with explicit assistance rules. Calculate true cash and normalized effort, not only tuition receipts and scheduled teaching hours.

\paragraph*{Week 8: Decide from evidence.} Repeat, narrow, change pricing/scope, train additional capacity, or stop. Preserve the actual learning from the work in a maintained operating record. State clearly what this small experiment does \textbf{not} prove about a later 6-month program or institution-wide scale.

\section*{Appendix H: The Operating Review and Capacity Decision Memorandum}

\begin{quote}
\small
\textbf{DATE / OFFER / COHORT OR CLIENT:}\\
\textbf{CURRENT STAGE:} exploratory $|$ 1st paid delivery $|$ repeated delivery $|$ delegated delivery\\[1ex]
\textbf{CUSTOMER VALUE:}\\
The precise promised result is\dots\\
The strongest outcome evidence is\dots\\
The main quality limitation is\dots\\[1ex]
\textbf{DEMAND:}\\
Who has paid, on what terms, from which channel?\\
Which future volume is observed or contracted versus forecast?\\[1ex]
\textbf{COST:}\\
Cash contribution under stated cost inclusion is\dots\\
Normalized contribution with founder replacement labor is\dots\\
Costs we still cannot estimate confidently are\dots\\[1ex]
\textbf{CAPACITY:}\\
Current binding constraint is\dots\\
At a 25\% increase in demand, the 1st failure would be\dots\\
The next capacity step would cost\dots\\[1ex]
\textbf{REUSABLE ASSETS:}\\
Which earlier investment demonstrably saved work or improved outcomes?\\
Which new asset is proposed, and why now?\\[1ex]
\textbf{AUTOMATION:}\\
The specific manual task is\dots\\
The quality and full-cost comparison shows\dots\\
The data/IP permissions are\dots\\[1ex]
\textbf{RISKS:}\\
Our financial affordable-loss limit is\dots\\
Our non-negotiable quality and privacy boundaries are\dots\\
Our fallback or stop condition is\dots\\[1ex]
\textbf{DECISION:}\\
Continue $|$ redesign $|$ hold enrollment $|$ invest $|$ delegate $|$ retire $|$ stop\\
Named owner, expected evidence, and review date\dots
\end{quote}
